\chapter{Introduction}

Bifurcation theory studies how the qualitative behaviour of a dynamical
system changes as a parameter is varied. This dissertation examines the
subcritical pitchfork bifurcation, a canonical example in which a stable
equilibrium loses stability discontinuously as unstable branches collide
with it, in contrast to the smooth, continuous loss of stability seen in
the supercritical case \citep{strogatz2015nonlinear}.

This is an example of a citation. \citet{taylor1923stability} is a classic
early example of a bifurcation analysis, applied to the stability of
Couette flow between rotating cylinders. Sensitivity to initial
conditions in nonlinear systems more generally was made famous by
\citet{lorenz1963deterministic}.

Figure~\ref{fig:subcritical_pitchfork_bifurcation} shows the bifurcation
diagram for the model system studied here. Table~\ref{tab:example} lists
a small set of computed equilibrium values used to validate the
numerical continuation code against the closed-form solution.

\begin{figure}[t]
\centering
\begin{subfigure}[b]{0.8\textwidth}
\centering
\resizebox{\linewidth}{!}{% This file was created by matlab2tikz.
% Minimal pgfplots version: 1.3
%
%The latest updates can be retrieved from
%  http://www.mathworks.com/matlabcentral/fileexchange/22022-matlab2tikz
%where you can also make suggestions and rate matlab2tikz.
%
\begin{tikzpicture}

\begin{axis}[%
width=0.95092\figurewidth,
height=\figureheight,
at={(0\figurewidth,0\figureheight)},
scale only axis,
separate axis lines,
every outer x axis line/.append style={black},
every x tick label/.append style={font=\color{black}},
xmin=-2,
xmax=2,
xlabel={$r$},
xmajorgrids,
every outer y axis line/.append style={black},
every y tick label/.append style={font=\color{black}},
ymin=-1.5,
ymax=1.5,
ylabel={$x$},
ymajorgrids,
legend style={legend cell align=left,align=left,draw=black}
]
\addplot [color=red,dashed,line width=2.0pt]
  table[row sep=crcr]{%
0.0202020202020203	0\\
0.0606060606060606	0\\
0.101010101010101	0\\
0.141414141414141	0\\
0.181818181818182	0\\
0.222222222222222	0\\
0.262626262626263	0\\
0.303030303030303	0\\
0.343434343434343	0\\
0.383838383838384	0\\
0.424242424242424	0\\
0.464646464646465	0\\
0.505050505050505	0\\
0.545454545454545	0\\
0.585858585858586	0\\
0.626262626262626	0\\
0.666666666666667	0\\
0.707070707070707	0\\
0.747474747474747	0\\
0.787878787878788	0\\
0.828282828282828	0\\
0.868686868686869	0\\
0.909090909090909	0\\
0.949494949494949	0\\
0.98989898989899	0\\
1.03030303030303	0\\
1.07070707070707	0\\
1.11111111111111	0\\
1.15151515151515	0\\
1.19191919191919	0\\
1.23232323232323	0\\
1.27272727272727	0\\
1.31313131313131	0\\
1.35353535353535	0\\
1.39393939393939	0\\
1.43434343434343	0\\
1.47474747474747	0\\
1.51515151515152	0\\
1.55555555555556	0\\
1.5959595959596	0\\
1.63636363636364	0\\
1.67676767676768	0\\
1.71717171717172	0\\
1.75757575757576	0\\
1.7979797979798	0\\
1.83838383838384	0\\
1.87878787878788	0\\
1.91919191919192	0\\
1.95959595959596	0\\
2	0\\
};
\addlegendentry{unstable};

\addplot [color=green,solid,line width=2.0pt]
  table[row sep=crcr]{%
-2	-0\\
-1.95959595959596	-0\\
-1.91919191919192	-0\\
-1.87878787878788	-0\\
-1.83838383838384	-0\\
-1.7979797979798	-0\\
-1.75757575757576	-0\\
-1.71717171717172	-0\\
-1.67676767676768	-0\\
-1.63636363636364	-0\\
-1.5959595959596	-0\\
-1.55555555555556	-0\\
-1.51515151515152	-0\\
-1.47474747474747	-0\\
-1.43434343434343	-0\\
-1.39393939393939	-0\\
-1.35353535353535	-0\\
-1.31313131313131	-0\\
-1.27272727272727	-0\\
-1.23232323232323	-0\\
-1.19191919191919	-0\\
-1.15151515151515	-0\\
-1.11111111111111	-0\\
-1.07070707070707	-0\\
-1.03030303030303	-0\\
-0.98989898989899	-0\\
-0.949494949494949	-0\\
-0.909090909090909	-0\\
-0.868686868686869	-0\\
-0.828282828282828	-0\\
-0.787878787878788	-0\\
-0.747474747474747	-0\\
-0.707070707070707	-0\\
-0.666666666666667	-0\\
-0.626262626262626	-0\\
-0.585858585858586	-0\\
-0.545454545454545	-0\\
-0.505050505050505	-0\\
-0.464646464646465	-0\\
-0.424242424242424	-0\\
-0.383838383838384	-0\\
-0.343434343434343	-0\\
-0.303030303030303	-0\\
-0.262626262626263	-0\\
-0.222222222222222	-0\\
-0.181818181818182	-0\\
-0.141414141414141	-0\\
-0.101010101010101	-0\\
-0.0606060606060606	-0\\
-0.0202020202020201	-0\\
0.0202020202020203	0\\
};
\addlegendentry{stable};

\addplot [color=red,dashed,line width=2.0pt,forget plot]
  table[row sep=crcr]{%
-2	1.4142135623731\\
-1.95959595959596	1.39985569241831\\
-1.91919191919192	1.38534902432272\\
-1.87878787878788	1.37068883368468\\
-1.83838383838384	1.35587014067861\\
-1.7979797979798	1.34088769029319\\
-1.75757575757576	1.32573593055923\\
-1.71717171717172	1.31040898851149\\
-1.67676767676768	1.29490064358918\\
-1.63636363636364	1.27920429813366\\
-1.5959595959596	1.26331294458641\\
-1.55555555555556	1.24721912892465\\
-1.51515151515152	1.23091490979333\\
-1.47474747474747	1.21439181269781\\
-1.43434343434343	1.19764077850724\\
-1.39393939393939	1.18065210538049\\
-1.35353535353535	1.16341538305773\\
-1.31313131313131	1.14591941825388\\
-1.27272727272727	1.12815214963553\\
-1.23232323232323	1.11010055054632\\
-1.19191919191919	1.09175051725163\\
-1.15151515151515	1.07308673997732\\
-1.11111111111111	1.05409255338946\\
-1.07070707070707	1.03474976236145\\
-1.03030303030303	1.0150384378451\\
-0.98989898989899	0.994936676326182\\
-0.949494949494949	0.974420314594759\\
-0.909090909090909	0.953462589245592\\
-0.868686868686869	0.932033727226042\\
-0.828282828282828	0.910100449556437\\
-0.787878787878788	0.887625364598595\\
-0.747474747474747	0.864566219253764\\
-0.707070707070707	0.840874965182522\\
-0.666666666666667	0.816496580927726\\
-0.626262626262626	0.791367567103066\\
-0.585858585858586	0.765413996382733\\
-0.545454545454545	0.738548945875996\\
-0.505050505050505	0.710669054518701\\
-0.464646464646465	0.681649810860727\\
-0.424242424242424	0.65133894727893\\
-0.383838383838384	0.619546918189723\\
-0.343434343434343	0.586032715327688\\
-0.303030303030303	0.55048188256318\\
-0.262626262626263	0.512470743190538\\
-0.222222222222222	0.471404520791032\\
-0.181818181818182	0.426401432711221\\
-0.141414141414141	0.376050716545178\\
-0.101010101010101	0.317820863081864\\
-0.0606060606060606	0.246182981958665\\
-0.0202020202020201	0.14213381090374\\
0.0202020202020203	0\\
};
\addplot [color=red,dashed,line width=2.0pt,forget plot]
  table[row sep=crcr]{%
-2	-1.4142135623731\\
-1.95959595959596	-1.39985569241831\\
-1.91919191919192	-1.38534902432272\\
-1.87878787878788	-1.37068883368468\\
-1.83838383838384	-1.35587014067861\\
-1.7979797979798	-1.34088769029319\\
-1.75757575757576	-1.32573593055923\\
-1.71717171717172	-1.31040898851149\\
-1.67676767676768	-1.29490064358918\\
-1.63636363636364	-1.27920429813366\\
-1.5959595959596	-1.26331294458641\\
-1.55555555555556	-1.24721912892465\\
-1.51515151515152	-1.23091490979333\\
-1.47474747474747	-1.21439181269781\\
-1.43434343434343	-1.19764077850724\\
-1.39393939393939	-1.18065210538049\\
-1.35353535353535	-1.16341538305773\\
-1.31313131313131	-1.14591941825388\\
-1.27272727272727	-1.12815214963553\\
-1.23232323232323	-1.11010055054632\\
-1.19191919191919	-1.09175051725163\\
-1.15151515151515	-1.07308673997732\\
-1.11111111111111	-1.05409255338946\\
-1.07070707070707	-1.03474976236145\\
-1.03030303030303	-1.0150384378451\\
-0.98989898989899	-0.994936676326182\\
-0.949494949494949	-0.974420314594759\\
-0.909090909090909	-0.953462589245592\\
-0.868686868686869	-0.932033727226042\\
-0.828282828282828	-0.910100449556437\\
-0.787878787878788	-0.887625364598595\\
-0.747474747474747	-0.864566219253764\\
-0.707070707070707	-0.840874965182522\\
-0.666666666666667	-0.816496580927726\\
-0.626262626262626	-0.791367567103066\\
-0.585858585858586	-0.765413996382733\\
-0.545454545454545	-0.738548945875996\\
-0.505050505050505	-0.710669054518701\\
-0.464646464646465	-0.681649810860727\\
-0.424242424242424	-0.65133894727893\\
-0.383838383838384	-0.619546918189723\\
-0.343434343434343	-0.586032715327688\\
-0.303030303030303	-0.55048188256318\\
-0.262626262626263	-0.512470743190538\\
-0.222222222222222	-0.471404520791032\\
-0.181818181818182	-0.426401432711221\\
-0.141414141414141	-0.376050716545178\\
-0.101010101010101	-0.317820863081864\\
-0.0606060606060606	-0.246182981958665\\
-0.0202020202020201	-0.14213381090374\\
0.0202020202020203	-0\\
};
\end{axis}
\end{tikzpicture}%}
\end{subfigure}
\caption{Bifurcation diagram for the subcritical pitchfork bifurcation, $\dot{x} = rx + x^3$, where for $r < 0$ there is a stable equilibrium at $x = 0$ and two unstable equilibria at $x = \pm\sqrt{-r}$, and for $r > 0$ the equilibrium at $x = 0$ is unstable.}
\label{fig:subcritical_pitchfork_bifurcation}
\end{figure}

\begin{table}[hp]
\centering
\hrule \vspace{0.5em}
\begin{tabular}{cccccc}
$r$ & $-4$ & $-3$ & $-2$ & $-1$ & $0$ \\
$x^*=\sqrt{-r}$ & $2.000$ & $1.732$ & $1.414$ & $1.000$ & $0.000$ \\
\end{tabular}
\caption{Non-zero equilibrium branch $x^*=\sqrt{-r}$ of $\dot{x}=rx+x^3$ for a range of the bifurcation parameter $r$, used to check the numerical continuation code against the closed-form root.}
\vspace{0.5em} \hrule
\label{tab:example}
\end{table}

\section{Structure of this dissertation}

Chapter~\ref{ch:methods} sets out the numerical continuation method used
to trace the equilibrium branches shown above, and describes how the
stability of each branch was determined from the sign of the linearised
system's eigenvalue. Chapter~\ref{ch:conclusion} summarises the findings
and suggests directions for extending the analysis to a two-parameter
unfolding of the bifurcation.
